% TOP_PROBLEM: {"id": 30002225, "title": "Trace-Cone Classification after Razak–Jacelon Stabilization", "category_id": 4, "category": {"id": 4, "name": "algebra", "display_name": "Algebra", "description": "Group theory, ring theory, field theory, and algebraic structures.", "slug": "algebra", "order_index": 4, "created_at": "2026-07-31T15:26:25.671Z"}, "status": "open"}
% FIRST_POSTED: 2026-10-06
% ATTEMPT_STATUS: solved
% OPENAI_CLAIM: {"schema_version":1,"source_repo":"https://github.com/openai/math","source_commit":"adc7f1241b42e322a6451854ab7e4b4c146bf78a","release_date":"2026-10-06","match":"full","resolution":"proved","independently_reviewed":false,"second_pass":"agreed","families":[{"family":"301","title":"Trace cones and Razak–Jacelon stabilization","subject":"Operator algebras","manuscripts":[{"title":"The trace cone classifies Razak–Jacelon stabilizations","pdf_path":"preprints/The-trace-cone-classifies-Razak-Jacelon-stabilizations-September-25-2026/The-trace-cone-classifies-Razak-Jacelon-stabilizations-September-25-2026.pdf","readme_path":"preprints/The-trace-cone-classifies-Razak-Jacelon-stabilizations-September-25-2026/README.md","date":"2026-09-25","theorem_ref":"thm:main (sections/01-introduction.tex:44); pt:intertwining (sections/08-point-intertwining.tex:403)","statement_sources":[{"path":"preprints/The-trace-cone-classifies-Razak-Jacelon-stabilizations-September-25-2026/build/sections/01-introduction.tex","line":44,"label":"thm:main"},{"path":"preprints/The-trace-cone-classifies-Razak-Jacelon-stabilizations-September-25-2026/build/sections/08-point-intertwining.tex","line":403,"label":"pt:intertwining"}]}],"lean":null,"reasoning_trace":null}]}
\documentclass[11pt]{article}
\usepackage[margin=25mm]{geometry}
\usepackage{fontspec}
\setmainfont{Latin Modern Roman}
\usepackage{amsmath,amssymb,mathtools}
\usepackage{unicode-math}
\setmathfont{Latin Modern Math}
\usepackage{enumitem}
\usepackage{xurl,hyperref,bookmark}
\hypersetup{colorlinks=true,urlcolor=blue}
\setcounter{secnumdepth}{0}
\setlength{\parindent}{0pt}
\setlength{\parskip}{5pt}
\setlength{\emergencystretch}{3em}
\begin{document}
\section*{Trace-Cone Classification after Razak–Jacelon Stabilization}

\subsection{OpenAI's claim}
The full topological cone of extended lower-semicontinuous traces is claimed to classify separable nuclear algebras after Razak--Jacelon and compact stabilization, realizing prescribed cone maps despite arbitrary ideals and nondense finite domains. The cited manuscript argues for that result; its relation to this problem is specified below.

\subsection{Mathematical statement of the claimed result}
\paragraph{Mathematical claim 1.} For separable nuclear complex $C^*$-algebras $A,B$, let $T(A),T(B)$ be their cones of lower-semicontinuous tracial weights, including zero and infinite values. Use the topology $\tau_i\to\tau$ when $\limsup_i\tau_i((a-\varepsilon)_+)\le\tau(a)\le\liminf_i\tau_i(a)$ for every positive $a$ and $\varepsilon>0$. OpenAI claims
\[
T(A)\simeq T(B)\text{ as topological additive cones}\quad\Longrightarrow\quad
A\otimes\mathcal W\otimes\mathcal K\simeq B\otimes\mathcal W\otimes\mathcal K.
\]
Cone isomorphisms preserve addition, zero and positive scalar multiplication; $\mathcal W$ denotes the Razak--Jacelon algebra; $\mathcal K$ denotes compact operators. Tensor products are spatial. The final intertwining theorem also claims an isomorphism implementing a prescribed compatible trace-cone map and Cuntz-semigroup map on these stabilizations.

\subsection{Outline of the argument}
Manuscript 1 (family 301, 2026-09-25): Local rank controllers define error ideals and small stabilizing homomorphisms, allowing countably many approximation errors to be absorbed. Completely positive cone maps match trace moments while lower ideal comparisons preserve exactly where weights become infinite. Separate existence and norm-uniqueness steps then support approximate intertwining, producing an isomorphism with the prescribed trace-cone action.

\subsection{Scope relative to this problem}
Manuscript 1 (family 301, 2026-09-25): The main theorem classifies the stated stabilizations of all separable nuclear algebras by the full extended lower-semicontinuous trace cone, with arbitrary ideals; the converse is the canonical trace-cone invariance under these stabilizations.

\subsection{Formal verification}
Family 301 has no Lean scope page in this snapshot; no formalization of these manuscript results is listed.

This site has not executed the Lean files or checked their dependencies. These descriptions report the scope of the linked declarations, rather than a new verification of the proofs.

\subsection{Status caveat}
OpenAI attributes these results to an internal model. Its release notice cautions that results lacking full formalization can contain mistakes. This site has not independently assessed the proofs; a solved label records the scope of the reported claim.

\subsection{Sources}
\begin{itemize}
\item OpenAI, \emph{The trace cone classifies Razak–Jacelon stabilizations}, 2026-09-25; citation key \texttt{OAI:The-trace-cone-classifies-Razak-Jacelon-stabilizations-September-25-2026}. \href{https://github.com/openai/math/blob/adc7f1241b42e322a6451854ab7e4b4c146bf78a/preprints/The-trace-cone-classifies-Razak-Jacelon-stabilizations-September-25-2026/The-trace-cone-classifies-Razak-Jacelon-stabilizations-September-25-2026.pdf}{PDF (GitHub)}; \href{https://github.com/openai/math/raw/adc7f1241b42e322a6451854ab7e4b4c146bf78a/preprints/The-trace-cone-classifies-Razak-Jacelon-stabilizations-September-25-2026/The-trace-cone-classifies-Razak-Jacelon-stabilizations-September-25-2026.pdf}{PDF (direct)}; \href{https://github.com/openai/math/blob/adc7f1241b42e322a6451854ab7e4b4c146bf78a/preprints/The-trace-cone-classifies-Razak-Jacelon-stabilizations-September-25-2026/README.md}{Manuscript page and citation}; \href{https://github.com/openai/math/blob/main/preprints/The-trace-cone-classifies-Razak-Jacelon-stabilizations-September-25-2026/The-trace-cone-classifies-Razak-Jacelon-stabilizations-September-25-2026.pdf}{Latest version}; Statement source: \href{https://github.com/openai/math/blob/adc7f1241b42e322a6451854ab7e4b4c146bf78a/preprints/The-trace-cone-classifies-Razak-Jacelon-stabilizations-September-25-2026/build/sections/01-introduction.tex\#L44}{thm:main}; Statement source: \href{https://github.com/openai/math/blob/adc7f1241b42e322a6451854ab7e4b4c146bf78a/preprints/The-trace-cone-classifies-Razak-Jacelon-stabilizations-September-25-2026/build/sections/08-point-intertwining.tex\#L403}{pt:intertwining}.
\item \href{https://github.com/openai/math/blob/adc7f1241b42e322a6451854ab7e4b4c146bf78a/CONTENTS.md}{OpenAI catalogue entry}.
\end{itemize}
\end{document}
